\section{Modules and Long-Range Interactions}
\label{sec:intro-modules}

%%%%%%%%%%%%%%%%%%%%%%%%%%%%%%%%%%%%%%%%%%%%%%%%%%%%%%%%%%%%%%%%%%%%%%%%%%%%%
%% Subsection 0: Motivation, why modules?

Unfortunately, the model \textsf{2D CCNTC} has several shortcomings
which are highlighted by the constant-depth communication in the
previous section. First and foremost, it is difficult to realize physically. 
A single, contiguous
2D lattice that contains an entire quantum architecture
may be prohibitively large to manufacture or too difficult to control
coherently. In practice,
scalable experiments already prefer to use many
smaller quantum computers that communicate by means of shared
entanglement \cite{Monroe2012}. Second, for highly parallellized
algorithms, the amount of quantum information undergoing active
computation is so large that communicating data across the
quantum processor consumes more resources than the computation
itself. Finally, as quantum architects we are interested in distilling
the essential locality of a computation, and distinguishing it from
communication (which itself may be an interesting problem).

Running an
algorithm on a single \textsf{2D CCNTC} architecture makes it
sensitive to the geometry of the different computation sites and how they
are arranged on the same lattice. Calculating the circuit resources of different
geometric arrangements will have vastly different numerical constants but will asymptotically
be the same. We would like a model that abstracts away this sensitivity.
The key to characterizing this essential locality is the long-range interaction. Rather than simulating them using short-range interactions,
we account for them separately. We will define a
\emph{hybrid architectural model} as one which accounts for long-range
and short-range interactions separately and also constrains them
according to independent graphs.

In fact, an example of this hybrid approach has already been proposed
in the \emph{quantum multicomputer} of Van Meter et al. \cite{VanMeter2008}.
This hybrid model is also called a distributed quantum computer since it contains several communicating nodes operating in parallel.
We will abbreviate is \textsf{QMC}.
Our approach is
very similar in that we divide a computation, which normally occurs on a single connected qubit graph $G$, into smaller subdivisions connected in
a graph $\overline{G}$. If two qubits $\ket{u}$ and $\ket{v}$ are connected
in $G$, then their two containing subdivisions $\overline{u} \ni \ket{u}$
and $\overline{v} \ni \ket{v}$ are connected in $\overline{G}$.
Like \textsf{QMC}, we also consider ``flying qubits,'' or shared
entanglement to allow for long-range, inter-node interactions. Moreover,
we also
consider a constrained topology for the nodes, one where the connectivity
graph has bounded degree. Finally, we treat short-range interactions between
qubits in the same node as being faster and cheaper.

However, we divide our
circuit into \emph{modules} to emphasize the following differences from
\textsf{QMC} nodes.\footnote{This should not be
confused with the word ``modular'' as in ``modular arithmetic'' or as
referring to the modulus $m$ that we are trying to factor.}
Each module is a self-contained \textsf{2D CCNTC} lattice, where the qubits
are constrained to have nearest-neighbor interactions. Qubits in a node,
in contrast, are \textsf{AC} circuits with arbitrary connectivity. The
Van Meter paper is more empirical, providing numerical calculations
based on estimates of current physical technology as well as more realistic
circuit implementations of long-range teleportations. Our work is more
abstract, providing asymptotic formulae for circuit resources. Finally,
while both nodes and modules have configurable size, we have come to the
conclusion that modules should contain $O(n)$ qubits each, whereas
the Van Meter paper concludes that nodes have an optimal $O(1)$ size for 
factoring.
We discuss module size in more detail in Section \ref{subsec:module-size}.

A related work by Beals et al. \cite{Beals2012} describes how to achieve
low-depth, generic simulations of any \textsf{AC} circuit, either by
rearranging qubits using a sorting network on \textsf{NTC} or
by mapping it to \textsf{QMC}
(even with fixed-size nodes connected in a \textsf{LNN} topology). The
depth overhead in either case for an $n$-qubit \textsf{AC} circuit is $O(\log^2 n)$. Their first approach of \textsf{NTC} sorting networks can be used
as an alternative to hybrid architectures. As we will see later in
Chapter \ref{chap:factor-polylog}, the asymptotic depth of the generic sorting network approach will be the same as our ``hand-crafted'' architecture for factoring,
but possibly with worse constants and asymptotically larger width and size. 

We have now described the motivation for a hybrid architectural model which
allows us to capture the essential locality of a quantum algorithm while
still accounting for long-range interactions. Consequently, we can
extend the model \textsf{2D CCNTC} to a further submodel which includes
modules.

%%%%%%%%%%%%%%%%%%%%%%%%%%%%%%%%%%%%%%%%%%%%%%%%%%%%%%%%%%%%%%%%%%%%%%%%%%%%
\subsection{The Submodel \textsf{2D CCNTCM}}
\label{subsec:2dccntcm}

The main feature of our new submodel is that modules
and long-range teleportations between them are nodes and edges, respectively,
in a higher-level planar graph. The teleportations each transmit one qubit
from one module to another, from any location within the source module
to any location within the destination module, making use of the
omnipresent classical controller. The modules can be arbitrarily far
apart physically, but they have bounded-degree (at most 4) connectivity with other
modules, and their edges are planar (they cannot intersect).

We do not discuss the physical technology underlying the long-range
teleportation other than acknowledging that it is based on shared
entanglement. Namely, the two modules which are the source and destination
of any long-range teleportation must each possess one half of an
Einstein-Podolsky-Rosen (EPR) pair, as shown in Equation \ref{eqn:epr}.

\begin{equation}
\normtwo \left( \ket{00} + \ket{11} \right)
\end{equation}

Therefore, each long-range teleportation includes a preprocessing step in
which the EPR pairs are generated and transported, along the bounded-degree
module graph, to the modules where future teleportations will take place.

A single module can be part of multiple teleportation operations in a single timestep, as long as they involve disjoint qubits within the module.
We justify this assumption in that it is
possible for a quantum computer to share multiple EPR pairs with a
second quantum computer, and even multiple other quantum computers,
simultaneously.

We call this new submodel \textsf{2D CCNTCM}, since it 
and we argue that is captures the essential aspects of 2D architectures
without being overly sensitive to the exact geometry of the lattices involved.
A graphic depiction of three modules in \textsf{2D CCNTCM} is shown in
Figure \ref{fig:modules}. Each module contains within it a
\textsf{2D CCNTC} lattice. We can equivalently consider the omnipresent,
single
classical controller as a collection of multiple classical controllers, one
for each module or teleportation operation, which can inter-communicate
classically and share a clock.

\textsf{2D CCNTCM} is very similar to the model of a quantum multicomputer
of Van Meter
\cite{VanMeter2006} with the difference that we allow parallel teleportations.
It is also very similar to the distributed quantum computer \textsf{DQC} of
Beals et al. \cite{Beals2012} except that we require each module to have
$\Omega(n)$ qubits in order to enforce some locality.

\begin{figure}[btp!]
\begin{center}
\includegraphics[width=4in]{factor-polylog/figures/modules.pdf}
\end{center}
\caption{Three modules in the \textsf{2D CCNTCM} model}
\label{fig:modules}
\end{figure}



The main point should be shifted, in the intro, and maybe elucidated here,
that low-depth for factoring can be achieved by moving into 2D, but
moreover that the width and size blowups can also be reduced with a
hybrid 2D module.

The model is hybrid, because instead of a single, contiguous 2D lattice,
say with $W$ qubits and $S$ interactions between qubits (counting the
general case of all single- and two-qubit gates as interactions),
the computation is split among many communicating 2D lattices
(say $\overline{W}$ of them), each containing $W/\overline{W}$ qubits
each. In this new model, we argue that there are now a total number of
qubit interactions $S'$ across all modules, where $S' \le S$, and at
the same time we have introduced $\overline{S}$
new, \emph{long-range} interactions between modules. Based on experimental
data of create shared entanglement between separate ion traps (citation needed here),
these long-range interactions are more expensive than short-range interactions
within a module, and so we count them separately.

Currently, we impose no constraints on the connectivity of the modules, which exist
on a higher-level graph of modules. This is supported by current proposals
(cite MUSIQC here). However, it is interesting to see what effect such constraints
would have on a factoring architecture. We will return to this later.


With no such constraints (arbitrary connectivity of modules), we can now see the
effect of setting the size of each module. If we set modules
to be a single qubit in size ($\overline{W} = W$), we would get the architectural
model \textsf{AC}. If we were to set module size to include the entire lattice
($\overline{W} = 1$), then we would get the model described in the previous
section, \textsf{2D CCNTC}. This new model then, \textsf{2D CCNTCM}, is a hybrid
between \textsf{AC} and \textsf{2D CCNTC} with module size as a parameter.
\textsf{AC} and \textsf{2D CCNTC} represent the two extremes of long-range
interactions, where \textsf{AC} permits all long-range interactions and
\textsf{2D CCNTC} permits none.

There are now two more things to study and discuss. One is the optimal setting
of the module size $W/\overline{W}$, which in this paper we have set to
be $O(n)$. The second is the connectivity constraints of the modules.

For the first parameter, we have not rigorously shown that the optimal setting
of $W / \overline{W}$ is 

Now then, what is the effect of imposing constraints on module connectivity?



\begin{definition}
A \textsf{2D CCNTCM} architecture consists of

\begin{itemize}
\item a quantum computer $\overline{QC}$ which is represented by a planar graph $(\overline{V},\overline{E})$. A
node $\overline{v} \in \overline{V}$ represents a module, or a graph $(V,E)$
from a \textsf{2D CCNTC} architecture defined previously. It can have
unbounded degree.
An
undirected edge $(\overline{u},\overline{v}) \in \overline{E}$ represents an
allowed teleportation from any qubit in module $\overline{u}$ to
another qubit in module $\overline{v}$.
\item All modules are restricted to be linear in the number of their qubits:
$|V| = \Theta(n)$ for all $(V,E) \in \overline{V}$.
\item a universal gate set $\mathcal{G} = \{X, Z, H, T, T^{\dagger}, CNOT,
MeasureZ\}$
for the qubits \emph{within the same} modules which is the same as for \textsf{2D CCNTC},
and an additional operation $Teleport$ which only operates on qubits
\emph{in
different} modules.
\item a deterministic machine (classical controller) $\overline{CC}$ that applies a sequence
of concurrent gates in each of $D+\overline{D}$ timesteps.
This can be a separate classical controller
for every pair of modules.
\item In timestep $i$, $\overline{CC}$ applies
gates $G_i = \{g_{i,j} : g_{i,j} \in \mathcal{G} \lor g_{i,j} = Teleport \}$.
That is, there are two kinds of timesteps with respect to the kinds of gates
which operate within them.
\begin{enumerate}
\item In the first kind, gates are exclusively from $\mathcal{G}$, and
they operate within modules as described
for \textsf{2D CCNTC} above. We say there are $D$ such timesteps.
\item In the second kind, gates are exclusively $Teleport$ gates between two qubits $v^{(1)}_{i,j} \in \overline{v}_1$ and
$v^{(2)}_{i,j} \in \overline{v}_2$ for
(possibly non-distinct) modules $\overline{v}_1, \overline{v}_2 \in \overline{V}$.
Again, all such qubits much be distinct within a timestep.
We say there are $\overline{D}$ such timesteps.
\end{enumerate}

Again, we define the support of $G_i$
as $V_i$, the set of all qubits acted upon by any $g_{i.j}$, which
includes all the modules.
\begin{equation}
V_i = \bigcup_{j: g_{i,j} \in G_i} v_{i,j} \cup v^{(1)}_{i,j} \cup v^{(2)}_{i,j} 
\end{equation}

\end{itemize}
\end{definition}

We measure the efficiency of a circuit in this new module using not just
the three conventional circuit resources, but with three novel resources
based on modules.

\begin{description}

%, depicted in Figure \ref{fig:resources}:
\item[module depth ($\overline{D}$):] the depth of consecutive teleportations between modules.
\item[module size ($\overline{S}$):] the number of total qubits teleported between any two modules over all timesteps.
\item[module width ($\overline{W}$):] the number of modules whose qubits are
acted upon during any timestep.

\end{description}

%We can make an observation analogous to Equation \ref{eqn:depth-width} but
%for modules in Equation \ref{eqn:module-depth-width}.

%\begin{equation}
%\overline{D} \le \overline{S} \le \overline{D}\cdot \overline{W}
%\label{eqn:module-depth-width}
%\end{equation}

We note the following relationship between circuit width and
module width.

\begin{equation}
W = O(n\overline{W})
\label{eqn:module-width}
\end{equation}

This restriction imposes some locality on our model by constraining it to
nearest-neighbor gates within a linear-sized group of qubits, but allowing
it long-range teleportation to circumvent onerous geometric constraints.
Using the constant-depth communication in Section \ref{sec:cdc}, and for
the specific case of factoring, we
can simulate arbitrary connectivity between modules with only a polynomial
increase in the module size and a constant increase in module depth.
